\section{Introduction}
\label{sec:intro}

Defragmentation has a bad reputation on modern storage: moving whole files costs
writes, slows co-running applications and wears out flash. FragPicker~\cite{park2021fragpicker}
answered this by tracing what an application actually reads and writes and moving
only the pieces whose requests straddle extent boundaries. Its artifact is a chain of
Python and shell scripts: an \emph{analysis phase} that turns a BCC trace into
per-file lists of hot regions, and a \emph{migration phase} that rewrites them.

The analysis phase processes files one at a time. An earlier course project by the
author of this paper added a parallel mode to it~\cite{forkbranch}: one batch
inode-to-path map, extent lookups through the FIEMAP ioctl, and a pool of threads or
processes. It reported a 12.74$\times$ speedup with 16 threads when a 5\,ms delay was
added in user space before every per-file extent lookup, and a slowdown without the
delay. The delay stood in for a slow disk or a network file system; no slow device
was measured. This paper asks whether that result holds up, and finds that it does
not, for reasons that are instructive beyond this one tool.

We re-evaluate the analysis phase from scratch: we read both code bases, capture
real traces with the artifact's own tracer, build an ablation ladder from the
unmodified upstream code to the parallel mode with a single-threaded rewrite in the
middle, and replace the user-space delay with delays injected at the block layer.
Our findings:

\begin{itemize}
\item \textbf{The parallel mode parallelises work the original never does.} Upstream's
  analysis phase reads no extents; the parallel mode reads every traced file's
  extents and then discards them. Its only effect on the output is to drop files on
  file systems without FIEMAP, where it lists none at all (\S\ref{sec:correct}).
\item \textbf{The original's cost is process creation and repeated tree walks.}
  \textbf{[E2/E3 numbers: Phase 1 per line, Phase 2 per file, scaling with $T$.]}
\item \textbf{A single thread beats every parallel configuration.} A 200-line rewrite
  that produces byte-identical output is \textbf{[X]}$\times$ faster than upstream and
  \textbf{[Y]}$\times$ faster than the best parallel configuration on the same traces
  (\S\ref{sec:ladder}).
\item \textbf{The speedup came from the latency model.} With a 5\,ms user-space delay
  we reproduce a \textbf{[Z]}$\times$ speedup; with the same delay injected at the block
  layer, files created together need \textbf{[R]} device reads in total and threads
  gain \textbf{[S]}; on a device that serves one request at a time they gain nothing
  (\S\ref{sec:models}).
\item \textbf{The trace itself is lossy.} Upstream's stop procedure loses the last
  buffered chunk of events, positions above 4\,GiB wrap around, and inode numbers
  from other file systems can be attributed to files under \texttt{/mnt}
  (\S\ref{sec:correct}).
\end{itemize}

None of these problems needs new hardware or clever parallelism to fix; they need a
competent sequential baseline~\cite{mcsherry2015cost}, latency injected where real
latency occurs, and a count of the device reads the work actually requires. We
release all scripts, raw data and traces (Appendix~\ref{sec:artifact}).
